\documentclass[11pt,letterpaper]{article}
\usepackage[utf8]{inputenc}
\usepackage[T1]{fontenc}
\usepackage{geometry}
\geometry{margin=2cm}
\usepackage{hyperref}
\usepackage{setspace}
\usepackage{siunitx}
\sisetup{separate-uncertainty=true, range-phrase=--}
\DeclareSIUnit\kcalmol{kcal\,mol^{-1}}
\DeclareSIUnit\angstrom{\text{\AA}}
\singlespacing

\begin{document}
\emergencystretch=2em

\begin{flushleft}
\textbf{Myke Vital Sao Temgoua} \\
Department of Physics, Faculty of Science \\
University of Yaound\'{e} I \\
P.O. Box 812, Yaound\'{e}, Cameroon \\
\href{mailto:myke-vital.sao@facsciences-uy1.cm}{myke-vital.sao@facsciences-uy1.cm} \\[1.5em]
September 11, 2026
\end{flushleft}

\begin{flushleft}
Prof.\ Kenneth M.\ Merz, Jr., Editor-in-Chief \\
\textit{Journal of Chemical Information and Modeling} \\
American Chemical Society
\end{flushleft}

\vspace{0.5em}

\noindent Dear Prof.\ Merz,

We submit our manuscript \textbf{``Docking--Molecular Dynamics Divergence as a Discovery Signal for African Antimalarial Natural Products''} for consideration as an Article in \textit{JCIM}.

The central question is whether the docking--MD divergence rate constitutes a reliable diagnostic signal for African antimalarial natural products. Seventeen candidates from a library of \num{65856} molecules generated from 396 African natural products were ranked by a docking-derived resistance-retention score (RRS) across six PfDHFR and PfCRT resistance states, with chemical-space (ACSI) and network (PNS) scores kept as separate ranking dimensions. The primary analysis on 12 two-target candidates yielded a discriminant five-class stratification (A*: 1, A: 1, B: 4, C: 5, D: 1) that separates predicted wild-type potency from predicted mutant retention. A 16-system single-replicate MD pilot on two Set-C candidates revealed that docking-RRS and short-MD geometry metrics diverged in direction for seven of eight mutant comparisons (87.5\% divergence rate), demonstrating that the two quantities are non-equivalent. This elevated divergence rate suggests that African natural products---with their rigid polycyclic frameworks, macrocyclic systems, and rare functional groups---may be susceptible to static-pose artifacts, though whether the rate reflects properties specific to this chemical class or a more general artefact of the short-MD protocol cannot be resolved here. An external docking panel (39 ligands, 312 records) confirmed that docking-score retention is protocol-dependent. The study supports computational prioritization but does not establish biological target engagement, affinity, or resistance resilience.

Code, source data, and analysis scripts are available on GitHub under an open-source licence; a Zenodo snapshot will be deposited at publication. Statistics are reported with effect sizes, permutation $p$-values, bootstrap intervals, and Bonferroni adjustment.

\noindent\textbf{Suggested reviewers.}
\begin{enumerate}
  \item Dr.\ Fidele Ntie-Kang, University of Buea --- African natural-product computational discovery
  \item Dr.\ Kelly Chibale, University of Cape Town --- antimalarial drug discovery
  \item Dr.\ John D.\ Chodera, Memorial Sloan Kettering --- free-energy methods and MD best practices
\end{enumerate}

\noindent This manuscript is original, not under consideration elsewhere, and all authors have approved the submission.

\vspace{1em}

\noindent Sincerely,

\vspace{1em}

\noindent \textbf{Myke Vital Sao Temgoua} \\
(on behalf of all co-authors)

\end{document}
